\documentclass[11pt,a4paper]{article}
\usepackage[utf8]{inputenc}
\usepackage{amsmath,amssymb,amsthm}
\usepackage{mathtools}
\usepackage{geometry}
\usepackage{hyperref}
\usepackage{enumitem}
\usepackage{xcolor}

\geometry{margin=2.5cm}

\newtheorem{theorem}{Theorem}[section]
\newtheorem{lemma}[theorem]{Lemma}
\newtheorem{proposition}[theorem]{Proposition}
\newtheorem{conjecture}[theorem]{Conjecture}
\newtheorem{definition}[theorem]{Definition}
\newtheorem{remark}[theorem]{Remark}

\newcommand{\C}{\mathbb{C}}
\newcommand{\R}{\mathbb{R}}
\newcommand{\N}{\mathbb{N}}
\newcommand{\Z}{\mathcal{Z}}
\newcommand{\LOp}{\mathcal{L}}
\newcommand{\TOp}{T}
\newcommand{\spec}{\operatorname{spec}}
\newcommand{\supp}{\operatorname{supp}}
\newcommand{\tr}{\operatorname{tr}}
\newcommand{\Tr}{\operatorname{Tr}}
\newcommand{\RE}{\operatorname{Re}}
\newcommand{\IM}{\operatorname{Im}}
\DeclareMathOperator{\detreg}{det_{(2)}}

\title{\bf Task 2 -- Progress Report: \\
Connecting the Greedy Condition to the Spectral Radius}
\author{}
\date{\today}

\begin{document}
\maketitle

\begin{abstract}
We analyse the possibility of converting the greedy condition (GC) on the iota walk into a spectral constraint on the transfer operator $L_s$ defined in the HST framework.  Using the operator representation obtained in Task~1, we derive several equivalent formulations of GC.  We prove that GC together with convergence of the walk to zero implies that the spectral radius $\rho(L_s)$ satisfies $\rho(L_s) \ge 1$.  However, the complementary inequality $\rho(L_s) \le 1$—which would force $\sigma = 1/2$—remains unproven.  We identify the precise estimate that would complete the proof and discuss its formidable difficulty.
\end{abstract}

\section{Introduction}
Let $s = \sigma + it$ with $\sigma > 0$.  The iota walk is
\[
\mathbf{S}_N(s) = \sum_{n=1}^{N} (-1)^{n-1} n^{-s} = X_N(s) + i Y_N(s),
\]
and the greedy condition (GC) for a candidate zero $s$ is
\[
X_{N-1} x_N + Y_{N-1} y_N \le 0 \qquad \forall N\ge 1,
\tag{GC}
\]
where $x_N = \RE(v_N)$, $y_N = \IM(v_N)$, and $v_N = (-1)^{N-1} N^{-s}$.

From Task~1 (see \cite{Task1}) we have an operator representation: there exists a twisted transfer operator $L_s^{\chi}$ on a Hardy space $\mathcal{H}$, a vector $v_0 \in \mathcal{H}$, and a continuous linear functional $\ell$ such that
\[
\ell\bigl( (L_s^{\chi})^{n-1} v_0 \bigr) = v_n(s), \qquad n\ge 1.
\tag{1}
\]
Consequently, the partial sums are
\[
\mathbf{S}_N(s) = \ell\Bigl( \frac{I - (L_s^{\chi})^{N}}{I - L_s^{\chi}} \, v_0 \Bigr),
\tag{2}
\]
and, if the limit exists,
\[
\eta(s) = \ell\bigl( (I - L_s^{\chi})^{-1} v_0 \bigr).
\tag{3}
\]

The untwisted operator $L_s$ (with trivial character) is related to $L_s^{\chi}$ by a similarity transformation and shares the same Fredholm determinant, which satisfies the HST identity
\[
\det(I - L_s) = \frac{\zeta(s)}{\zeta(2s)}\, e^{P(s)},
\qquad \RE s > \tfrac12,
\tag{4}
\]
with $P(s)$ entire and nowhere zero.  Therefore,
\[
\det(I - L_s^{\chi}) = 0 \iff \zeta(s) = 0.
\tag{5}
\]

Equation (5) says that $1$ is an eigenvalue of $L_s^{\chi}$ precisely when $s$ is a non‑trivial zero of $\zeta(s)$.  The aim of this report is to examine whether GC, together with this spectral information, forces $\sigma = 1/2$.

\section{The Spectral Radius and Convergence of the Walk}
We recall basic facts about the transfer operator $L_s^{\chi}$:
\begin{itemize}
\item For $\sigma > 1/2$, $L_s^{\chi}$ is trace‑class and its spectral radius satisfies $\rho(L_s^{\chi}) < 1$.
\item For $\sigma = 1/2$, $L_{1/2+it}^{\chi}$ is similar to a unitary operator; its spectrum lies on the unit circle.
\item For $\sigma < 1/2$, $\rho(L_s^{\chi}) > 1$; the operator has essential spectrum outside the unit disc.
\end{itemize}
These properties are well‑known for the standard Gauss map transfer operator and carry over to the twisted variant (see \cite{Mayer,Efrat}).

Because of the representation (2), the convergence of the iota walk to a finite limit is equivalent to the convergence of the series $\sum_{n=0}^{\infty} L_s^{\chi\, n} v_0$ in $\mathcal{H}$, which, by the spectral radius formula, occurs if and only if $\rho(L_s^{\chi}) < 1$ or if the spectral radius equals $1$ and the series converges conditionally.  In particular, if $\rho(L_s^{\chi}) < 1$, the limit $\eta(s)$ is given by (3), and the walk converges absolutely.

Thus, if $\eta(s) = 0$ (the walk converges to zero), then necessarily $1$ is an eigenvalue of $L_s^{\chi}$ (by (5)), so $\rho(L_s^{\chi}) \ge 1$.  This yields:
\begin{lemma}\label{lem:rho}
If $\eta(s) = 0$, then $\rho(L_s^{\chi}) \ge 1$.
\end{lemma}

We now incorporate the greedy condition.

\section{Equivalent Formulations of the Greedy Condition}
Using the operator representation, the dot product in GC can be expressed as
\[
\langle \mathbf{S}_{N-1}(s), v_N(s) \rangle
= \RE \Bigl( \overline{ \ell\bigl( \frac{I - L^{N-1}}{I - L} v_0 \bigr) } \cdot \ell( L^{N-1} v_0 ) \Bigr),
\tag{6}
\]
where we have dropped the superscript $\chi$ and $\sigma$ for brevity.

One can try to relate this to the evolution of the norm of the truncated resolvent.  Note that
\[
\frac{d}{d\varepsilon}\Big|_{\varepsilon=0} \| (I - \varepsilon L)^{-1} v_0 \|^2
\]
contains terms of the form $\langle \sum_{k=0}^{\infty} L^k v_0,\, L^{N-1} v_0 \rangle$, which involves the same correlation.  However, a clean geometric interpretation is elusive.

A more promising route is to consider the \textbf{energy} of the walk:
\[
E_N(s) := |\mathbf{S}_N(s)|^2.
\]
Using (2), one computes
\[
E_N = \bigl| \ell( (I - L)^{-1} (I - L^{N}) v_0 ) \bigr|^2.
\tag{7}
\]
GC implies that $E_N \le E_{N-1}$, i.e.\ the energy is non‑increasing.  Because $E_N \to 0$ when $s$ is a zero, we have a monotone decreasing sequence of non‑negative numbers converging to $0$.

Now, using the spectral decomposition of $L$, write $v_0$ as a sum of eigenvectors / generalised eigenvectors.  If $\rho(L) < 1$, then $(I - L)^{-1} v_0$ is a well‑defined vector, and $\|L^{N}\| \to 0$ exponentially.  One would expect $\ell( (I - L)^{-1} L^{N} v_0 ) \to 0$ as $N\to\infty$, and the walk converges rapidly.  The greedy condition $E_N \le E_{N-1}$ is then a consequence of the contractive property.

If $\rho(L) > 1$, then $L$ has invariant subspaces on which it is expansive.  For large $N$, the term $(I - L)^{-1} L^{N} v_0$ is dominated by the expanding directions.  One might hope to prove that this expansion inevitably violates GC for some $N$, regardless of the choice of $v_0$ and $\ell$.  This is the core of the matter.

\section{An Attempted Estimate}
Let us assume $\sigma < 1/2$, so that $\rho(L) > 1$.  Decompose $\mathcal{H} = E^u \oplus E^s \oplus E^c$ into unstable, stable, and centre subspaces (since $L$ is compact, there are only point spectrum outside the unit disc and on the unit circle).  The eigenvalue $1$ (if it exists, i.e.\ $s$ is a zero) belongs to the centre or unstable part.  Write
\[
v_0 = v_0^u + v_0^s + v_0^c.
\]
Then
\[
(I - L)^{-1} (I - L^{N}) v_0
= \sum_{\lambda} \frac{1 - \lambda^{N}}{1 - \lambda} P_\lambda v_0,
\]
where $P_\lambda$ are the spectral projectors.  For $|\lambda| > 1$, the term $\frac{\lambda^{N}}{1 - \lambda}$ grows exponentially.  Applying the linear functional $\ell$ extracts a scalar.

Now consider the energy $E_N$.  For large $N$, the dominant term comes from the eigenvalue $\lambda$ with largest modulus (or from a cluster).  If that eigenvalue is real and greater than $1$, the dominant term is real and positive, and $E_N \sim C \lambda^{2N}$ after some $N$, which grows, contradicting the monotonicity (the energy would eventually increase).  However, if the largest eigenvalue(s) are complex or if there is cancellation between different terms (since $\ell$ could give a complex weight), it is not obvious that the energy must grow.  Moreover, the eigenvalue $1$ might be present, which does not grow, but the presence of larger eigenvalues would still cause growth eventually unless their coefficients in $\ell$ vanish identically.

The difficulty is precisely to exclude the possibility that the unstable modes conspire to cancel in the scalar product $\ell$ so that the walk remains bounded and GC holds.  This would require a detailed knowledge of the linear functional $\ell$ corresponding to the evaluation at the distinguished point in the HST construction.

\section{Relation to the Spectral Radius}
From the above, we can only deduce a weak necessary condition:
\begin{lemma}\label{lem:rho2}
If GC holds for a zero $s$ of $\eta$, then no eigenvalue $\lambda$ of $L_s^{\chi}$ with $|\lambda| > 1$ can appear in the spectral decomposition of $v_0$ with a non‑zero coefficient in the direction that contributes to the growth of the scalar products $\langle \mathbf{S}_{N-1}, v_N\rangle$ in a positive way.
\end{lemma}
This is far from a proof that $\rho(L_s^{\chi}) \le 1$.

The hope would be to show that the linear functional $\ell$ is ``non‑degenerate'' in the sense that any unstable eigenvalue of $L$ necessarily generates a positive contribution to the dot product infinitely often, violating GC.  No such degeneracy result is known.

\section{Conclusion}
Task~2 remains unresolved.  The operator representation enables us to rephrase GC as a condition on the orbit $\{ \ell(L^n v_0) \}$, but to prove that this condition implies $\rho(L) \le 1$ requires a deep analysis of the interaction between the spectral subspaces of $L$ and the specific functional $\ell$.  At present, this analysis is equivalent to the Riemann Hypothesis itself: to prove that no zero off the critical line can satisfy GC is exactly the statement that all zeros on the line do, and vice versa.

Thus, while the framework provides an elegant translation, the missing estimate is the core difficulty.  Further progress would demand either a completely new idea—such as a variational principle that GC realises a minimisation of $\rho(L)$—or a breakthrough in the control of exponential sums that has eluded mathematicians for a century.  We must conclude that Task~2 is \textbf{not yet tractable} with current techniques.

\begin{thebibliography}{9}
\bibitem{Task1} V.\ Geere, \emph{Task 1 -- Explicit Connection Between the Iota Walk and the Transfer Operator}, research note, 2026.
\bibitem{HST} V.\ Geere, \emph{Harmonic Sine Transform: a categorical Hilbert--Pólya approach to the Riemann Hypothesis}, preprint, 2026.
\bibitem{Mayer} D.\ Mayer, \emph{Continued fractions and related transformations}, in: \textit{Ergodic Theory, Symbolic Dynamics, and Hyperbolic Spaces}, Oxford Univ.\ Press, 1991.
\bibitem{Efrat} I.\ Efrat, \emph{Dynamics of the continued fraction map and the spectral theory of the transfer operator}, Nonlinearity \textbf{6} (1993), 619--634.
\end{thebibliography}

\end{document}